% GENERATED FILE. Edit canonical lessons, then run npm run generate:books.
% canonical-source-hash: fnv1a64:095ca583c72d12ea
% canonical-lessons: GE-C325-teig, GE-C325-hefe, GE-C325-kneten, GE-C325-ausrollen, GE-C325-nudelholz, GE-R325-review-318-325

\chapter{Der Teig}
\label{ch:ge-der-teig}
\hlchaptermodality{325}

\begin{chapteropening}
\textbf{By the end of this chapter:} \emph{I can describe making dough: dough, yeast, to knead, to roll out and a rolling pin.}

Complete the last lesson of chapter 325: I can describe making dough: dough, yeast, to knead, to roll out and a rolling pin.
\end{chapteropening}

\section[der Teig]{der Teig — the dough, the batter}

\label{lesson:GE-C325-teig}



\begin{quote}
Before the new one: say the German for to crack open, then the German for the whisk.
\end{quote}

\subsection*{You'll want to know: der Teig}
\textbf{der Teig} — \textquotedblleft{}the dough, the batter\textquotedblright{}.

\begin{quote}
Der Teig muss eine Stunde ruhen.
\end{quote}

\begin{quote}
\emph{The dough has to rest for an hour.}
\end{quote}

Plural: \emph{die Teige}.

\begin{cousinweb}
\emph{Teig} is inherited from Old High German \emph{teig}, from Proto-Germanic \emph{daigaz} (EtymDB), the same source as English \emph{dough}.
\end{cousinweb}

\begin{grammarlens}[title={What you've built}]
One more word for telling, explaining and describing at B1.
\end{grammarlens}

\subsection*{Your turn}
\begin{itemize}
\raggedright
  \item \emph{Say it:} \emph{der Teig}
  \item \emph{Say it:} \emph{der Teig}, once more
  \item \emph{Say it:} a sentence of your own with \emph{der Teig}
  \item \emph{Recall:} say the German for to crack open, then the German for the whisk, then say \emph{der Teig} again
\end{itemize}

\begin{tcolorbox}[breakable,colback=teal!4,colframe=teal!35!black,title={Before you move on}]
What does der Teig mean? (\textquotedblleft{}The dough, the batter\textquotedblright{}.) Say it once more.
\end{tcolorbox}
\section[die Hefe]{die Hefe — the yeast}

\label{lesson:GE-C325-hefe}



\begin{quote}
Before the new one: say the German for the whisk, then the German for the dough.
\end{quote}

\subsection*{You'll want to know: die Hefe}
\textbf{die Hefe} — \textquotedblleft{}the yeast\textquotedblright{}.

\begin{quote}
Für das Brot braucht man frische Hefe.
\end{quote}

\begin{quote}
\emph{For the bread you need fresh yeast.}
\end{quote}

No plural. In Austria: \emph{die Germ}.

\begin{cousinweb}
\emph{Hefe} is inherited from Old High German \emph{hevo}, from Proto-Germanic \emph{hafjô} \textquotedblleft{}yeast\textquotedblright{} (EtymDB), related to \emph{heben} \textquotedblleft{}to lift\textquotedblright{}: what makes dough rise.
\end{cousinweb}

\begin{grammarlens}[title={What you've built}]
One more word for telling, explaining and describing at B1.
\end{grammarlens}

\subsection*{Your turn}
\begin{itemize}
\raggedright
  \item \emph{Say it:} \emph{die Hefe}
  \item \emph{Say it:} \emph{die Hefe}, once more
  \item \emph{Say it:} a sentence of your own with \emph{die Hefe}
  \item \emph{Recall:} say the German for the whisk, then the German for the dough, then say \emph{die Hefe} again
\end{itemize}

\begin{tcolorbox}[breakable,colback=teal!4,colframe=teal!35!black,title={Before you move on}]
What does die Hefe mean? (\textquotedblleft{}The yeast\textquotedblright{}.) Say it once more.
\end{tcolorbox}
\section[kneten]{kneten — to knead}

\label{lesson:GE-C325-kneten}



\begin{quote}
Before the new one: say the German for the dough, then the German for the yeast.
\end{quote}

\subsection*{You'll want to know: kneten}
\textbf{kneten} — \textquotedblleft{}to knead\textquotedblright{}.

\begin{quote}
Er knetete den Teig zehn Minuten lang.
\end{quote}

\begin{quote}
\emph{He kneaded the dough for ten minutes.}
\end{quote}

Regular: \emph{knetete}, \emph{hat geknetet}.

\begin{cousinweb}
\emph{kneten} is inherited from Old High German \emph{knetan}, from Proto-Germanic \emph{knedaną} \textquotedblleft{}to knead\textquotedblright{} (EtymDB), the same source as English \emph{knead}.
\end{cousinweb}

\begin{grammarlens}[title={What you've built}]
One more word for telling, explaining and describing at B1.
\end{grammarlens}

\subsection*{Your turn}
\begin{itemize}
\raggedright
  \item \emph{Say it:} \emph{kneten}
  \item \emph{Say it:} \emph{kneten}, once more
  \item \emph{Say it:} a sentence of your own with \emph{kneten}
  \item \emph{Recall:} say the German for the dough, then the German for the yeast, then say \emph{kneten} again
\end{itemize}

\begin{tcolorbox}[breakable,colback=teal!4,colframe=teal!35!black,title={Before you move on}]
What does kneten mean? (\textquotedblleft{}To knead\textquotedblright{}.) Say it once more.
\end{tcolorbox}
\section[ausrollen]{ausrollen — to roll out}

\label{lesson:GE-C325-ausrollen}



\begin{quote}
Before the new one: say the German for the yeast, then the German for to knead.
\end{quote}

\subsection*{You'll want to know: ausrollen}
\textbf{ausrollen} — \textquotedblleft{}to roll out\textquotedblright{}.

\begin{quote}
Sie rollte den Teig dünn aus.
\end{quote}

\begin{quote}
\emph{She rolled the dough out thin.}
\end{quote}

Separable: \emph{sie rollt aus}, \emph{sie hat ausgerollt}.

\begin{cousinweb}
\emph{ausrollen} is built on \emph{rollen}, from \emph{Rolle} \textquotedblleft{}roll\textquotedblright{}, which came through Old French \emph{rolle} from Latin \emph{rotulus} \textquotedblleft{}a little wheel\textquotedblright{} (EtymDB).
\end{cousinweb}

\begin{grammarlens}[title={What you've built}]
One more word for telling, explaining and describing at B1.
\end{grammarlens}

\subsection*{Your turn}
\begin{itemize}
\raggedright
  \item \emph{Say it:} \emph{ausrollen}
  \item \emph{Say it:} \emph{ausrollen}, once more
  \item \emph{Say it:} a sentence of your own with \emph{ausrollen}
  \item \emph{Recall:} say the German for the yeast, then the German for to knead, then say \emph{ausrollen} again
\end{itemize}

\begin{tcolorbox}[breakable,colback=teal!4,colframe=teal!35!black,title={Before you move on}]
What does ausrollen mean? (\textquotedblleft{}To roll out\textquotedblright{}.) Say it once more.
\end{tcolorbox}
\section[das Nudelholz]{das Nudelholz — the rolling pin}

\label{lesson:GE-C325-nudelholz}



\begin{quote}
Before the new one: say the German for to knead, then the German for to roll out.
\end{quote}

\subsection*{You'll want to know: das Nudelholz}
\textbf{das Nudelholz} — \textquotedblleft{}the rolling pin\textquotedblright{}.

\begin{quote}
Mit dem Nudelholz rollte er den Teig aus.
\end{quote}

\begin{quote}
\emph{He rolled out the dough with the rolling pin.}
\end{quote}

Plural: \emph{die Nudelhölzer}. In Austria: \emph{der Nudelwalker}.

\begin{cousinweb}
\emph{Nudelholz} joins \emph{Nudel} \textquotedblleft{}noodle\textquotedblright{} and \emph{Holz} \textquotedblleft{}wood\textquotedblright{}, Old High German \emph{holz}, Proto-Germanic \emph{hultą} (EtymDB): the wooden tool for noodle dough.
\end{cousinweb}

\begin{grammarlens}[title={What you've built}]
One more word for telling, explaining and describing at B1.
\end{grammarlens}

\subsection*{Your turn}
\begin{itemize}
\raggedright
  \item \emph{Say it:} \emph{das Nudelholz}
  \item \emph{Say it:} \emph{das Nudelholz}, once more
  \item \emph{Say it:} a sentence of your own with \emph{das Nudelholz}
  \item \emph{Recall:} say the German for to knead, then the German for to roll out, then say \emph{das Nudelholz} again
\end{itemize}

\begin{tcolorbox}[breakable,colback=teal!4,colframe=teal!35!black,title={Before you move on}]
What does das Nudelholz mean? (\textquotedblleft{}The rolling pin\textquotedblright{}.) Say it once more.
\end{tcolorbox}
\section[(dialogue)]{Words from chapters 318–325 — a review}

\label{lesson:GE-R325-review-318-325}



\begin{quote}
Say the German for the war and the peace.
\end{quote}

\subsection*{Your turn}
Say these aloud:

\begin{itemize}
\raggedright
  \item the war, the peace, the battle, to defeat, the defeat
  \item the emperor, to rule, to govern, to conquer, the victory
  \item the Middle Ages, antiquity, the chronicle, the ancestor, the downfall
  \item to invent, the invention, to found, to come into being, to return
  \item the ingredient, the pinch, the tablespoon, to measure out, the bowl
  \item to peel, to chop, to dice, to grate, to halve
  \item to stir, to give a stir, to whisk, to crack open, the whisk
  \item the dough, the yeast, to knead, to roll out, the rolling pin
\end{itemize}

\begin{tcolorbox}[breakable,colback=teal!4,colframe=teal!35!black,title={Before you move on}]
The war? (\textbf{der Krieg}.) The peace? (\textbf{der Frieden}.) The battle? (\textbf{die Schlacht}.) To defeat? (\textbf{besiegen}.) The defeat? (\textbf{die Niederlage}.) The emperor? (\textbf{der Kaiser}.) To rule? (\textbf{herrschen}.) To govern? (\textbf{regieren}.) To conquer? (\textbf{erobern}.) The victory? (\textbf{der Sieg}.) The Middle Ages? (\textbf{das Mittelalter}.) Antiquity? (\textbf{die Antike}.) The chronicle? (\textbf{die Chronik}.) The ancestor? (\textbf{der Vorfahr}.) The downfall? (\textbf{der Untergang}.) To invent? (\textbf{erfinden}.) The invention? (\textbf{die Erfindung}.) To found? (\textbf{gründen}.) To come into being? (\textbf{entstehen}.) To return? (\textbf{zurückkehren}.) The ingredient? (\textbf{die Zutat}.) The pinch? (\textbf{die Prise}.) The tablespoon? (\textbf{der Esslöffel}.) To measure out? (\textbf{abmessen}.) The bowl? (\textbf{die Schüssel}.) To peel? (\textbf{schälen}.) To chop? (\textbf{hacken}.) To dice? (\textbf{würfeln}.) To grate? (\textbf{raspeln}.) To halve? (\textbf{halbieren}.) To stir? (\textbf{rühren}.) To give a stir? (\textbf{umrühren}.) To whisk? (\textbf{verquirlen}.) To crack open? (\textbf{aufschlagen}.) The whisk? (\textbf{der Schneebesen}.) The dough? (\textbf{der Teig}.) The yeast? (\textbf{die Hefe}.) To knead? (\textbf{kneten}.) To roll out? (\textbf{ausrollen}.) The rolling pin? (\textbf{das Nudelholz}.)
\end{tcolorbox}
